\chapter{Bugs between the services}
\label{ch:bugs-infra}

These are the hardest failures in the book. In every one, each service was
working perfectly and the system still did not.

\section{The dashboard that returned 200 and showed nothing}

This is the best example in the book, because every conventional diagnostic said
the system was healthy.

\textbf{Symptom.} \fpath{http://localhost:8080/horizon} rendered a blank dark
rectangle. No error message on the page.

\textbf{What all the usual checks said:}

\begin{center}
\begin{tabular}{@{}>{\raggedright\arraybackslash}p{6.6cm}p{7.0cm}@{}}
\toprule
\textbf{Check} & \textbf{Result} \\
\midrule
HTTP status of \fpath{/horizon} & \fpath{200 OK} \\
Server error logs & clean \\
Container health checks & all passing \\
The HTML shell & delivered correctly \\
\fpath{/horizon/api/stats} & \fpath{200}, valid JSON \\
\bottomrule
\end{tabular}
\end{center}

\textbf{Wrong guess \#1.} The dashboard's JavaScript is missing. Checking
\fpath{/vendor/horizon/app.js} returned \fpath{200} --- but with
\fpath{Content-Type: text/html} and a body beginning \fpath{<!doctype html>}.
That looked like the answer: nginx had no rule for \fpath{/vendor/}, so the
request fell through to the SPA, which returns \fpath{index.html} for unknown
paths. The browser asks for JavaScript and receives HTML.

That was a real routing flaw. It was not this bug.

\textbf{Wrong guess \#2.} Publish Horizon's assets to \fpath{public/vendor}.
Running the command printed: \emph{``Horizon no longer publishes its assets.''}
Modern Horizon serves them through its own routes, so \fpath{/vendor/horizon/}
was never the correct URL to begin with.

\textbf{The evidence that ended it.} Open the browser's developer console:

\begin{lstlisting}
Executing inline script violates the following Content Security Policy
directive 'script-src 'self''.

Loading the stylesheet 'https://fonts.bunny.net/css?family=figtree'
violates the following Content Security Policy directive: "style-src 'self'"
\end{lstlisting}

\textbf{The cause.} The site-wide Content Security Policy is deliberately
strict:

\where{infra/nginx/default.conf}
\begin{lstlisting}[language=nginxconf]
# No CDNs, no external fonts, no inline scripts. The SPA is entirely
# self-hosted, so the policy can be this tight -- and an incident dashboard
# should render even when the wider internet is the thing that is broken.
\end{lstlisting}

That reasoning is correct \emph{for the React app}. Horizon is a third-party
dashboard that boots from an inline \fpath{<script>}, loads a webfont from an
external host, and --- discovered one round later --- compiles Vue templates at
runtime, which requires \fpath{eval}. One policy could not serve both.

\subsection{The fix that did nothing, and why}

The obvious repair is a looser policy scoped to that one location:

\begin{lstlisting}[language=nginxconf]
location /horizon {
    add_header Content-Security-Policy "... 'unsafe-inline' ..." always;
    try_files $uri /index.php?$query_string;
}
\end{lstlisting}

It had no effect whatsoever. The response still carried the strict policy.

\note{Here is why, and it is a genuinely subtle piece of nginx behaviour.
\fpath{try\_files \$uri /index.php} performs an \textbf{internal redirect}: nginx
stops processing this location and starts again in the location that handles
\fpath{index.php}. Response headers come from whichever location \emph{finally}
handles the request --- so the \fpath{add\_header} in \fpath{/horizon} never
applied. Worse, \fpath{\$uri} is rewritten by that redirect and becomes
\fpath{/index.php}, so testing on \fpath{\$uri} would not work either.
\fpath{\$request\_uri} keeps the original path.}

\textbf{The fix.} Choose the policy with a \fpath{map} on \fpath{\$request\_uri},
and emit it once at server level.

\where{infra/nginx/default.conf}
\begin{lstlisting}[language=nginxconf]
map $request_uri $incidentflow_csp {
    default    "default-src 'self'; script-src 'self'; ...";
    ~^/horizon "default-src 'self'; script-src 'self' 'unsafe-inline' 'unsafe-eval'; ...";
}
\end{lstlisting}

\shot{14-horizon-dashboard.png}{Horizon after the fix. ``Jobs Per Minute'',
``Failed Jobs Past 7 Days'' --- the queue is finally visible.}

\warnbox{The relaxation is real, and the config says so rather than pretending
otherwise. \fpath{unsafe-inline} plus \fpath{unsafe-eval} is a meaningful
weakening. It is accepted because it is scoped to an authenticated operator
path, it cannot reach the application, and the alternative was having no queue
visibility at all. Write down what a compromise costs; a comment that only
records the benefit is marketing.}

\subsection{The lesson}

Every server-side signal said healthy. The only evidence existed in the
browser, in a console nobody had opened. \textbf{A class of failure exists that
your server can neither detect nor log}, because the rejection happens on the
other machine --- and a policy failure is exactly that class: nothing is broken,
everything behaves as specified, and the specification forbids the thing you
needed.

\section{The 502 that appeared while nothing changed}

\textbf{Symptom.} The application had been working. Later, every API request
through nginx returned \fpath{502 Bad Gateway}. No deploy, no code change.

\textbf{Wrong guess.} PHP-FPM has died or is listening on the wrong interface.
Checking inside the container showed it listening on \fpath{:::9000} and alive.

\textbf{Evidence.} The container uptimes:

\begin{lstlisting}
nginx    Up 4 hours
api      Up 3 hours
\end{lstlisting}

The API restarted \emph{after} nginx started.

\textbf{The cause.} nginx resolves the hostnames in an \fpath{upstream} block
\textbf{once, at startup}, and caches the address. The API container restarted
and Docker gave it a new IP. nginx was still dialling the old one.

\textbf{The fix.} \fpath{docker compose restart nginx} --- and understanding that
this will recur every time the API restarts alone.

\note{The same class of bug appears in the deployment configuration, where
upstreams are declared with a resolver and re-resolved at request time
specifically so a redeploy of one service does not leave the edge pointing at a
dead address. Recognising a bug you have already met, in a different disguise, is
most of what experience is.}

\section{The health check that lied}

\textbf{Symptom.} CI failed intermittently: \fpath{migrate didn't complete
successfully: exit 1}, while PostgreSQL reported \textbf{healthy}. Roughly one
run in two.

\textbf{Evidence.} Sampling both forms of the readiness probe on a fresh
database, every 0.3 seconds:

\begin{lstlisting}
t=00.0s   socket(no -h): --     tcp(-h 127.0.0.1): --
t=01.8s   socket(no -h): UP     tcp(-h 127.0.0.1): --      <-- healthy, nothing listening
t=02.1s   socket(no -h): --     tcp(-h 127.0.0.1): --
t=03.3s   socket(no -h): UP     tcp(-h 127.0.0.1): UP
\end{lstlisting}

\textbf{The cause.} The check was \fpath{pg\_isready -U incidentflow -d
incidentflow}. With no \fpath{-h}, \fpath{pg\_isready} connects over the
\textbf{Unix socket}. The official Postgres image runs its initialisation phase
on a temporary server started with \fpath{listen\_addresses=''} --- socket only,
no TCP. So the probe succeeds mid-initialisation, the container is marked
healthy, and everything waiting on it is released against a port that accepts
nothing.

There was even a comment claiming the check tested a real connection by naming
the user and database. It does not: \fpath{pg\_isready} never opens a database
connection, it asks the postmaster whether it is accepting them.

\textbf{The fix.} \fpath{pg\_isready -h 127.0.0.1 -U incidentflow -d
incidentflow}.

\warnbox{The general rule: \textbf{a readiness probe must use the same transport
as the client it gates.} Probing a Unix socket to certify TCP readiness tests a
different thing from the one that matters --- and the two diverge exactly during
startup, which is the only window the probe exists for.}

\section{Three smaller ones, quickly}

\subsection{A clean exit treated as a failure}

\fpath{docker compose up --wait} returned 1 with \fpath{migrate Exited (0)}. The
migration container does its job and stops, and \fpath{--wait} counted a stopped
container as a failure. The fix is \fpath{condition:
service\_completed\_successfully}, which also made \fpath{api}, \fpath{horizon}
and \fpath{scheduler} wait for the schema to exist --- something they had never
actually been told to do.

\subsection{Workers permanently unhealthy}

\fpath{horizon} and \fpath{scheduler} inherited the API image's health check,
which probes php-fpm on port 9000. Neither runs php-fpm. Both sat marked
unhealthy while working perfectly --- Horizon was visibly processing jobs while
reporting itself broken. Disabling an inherited check that describes a different
process is the fix.

\subsection{The second `up' that destroyed the stack}

Running \fpath{docker compose up} a second time, without removing volumes,
aborted everything: the seeder re-inserted blindly and hit a unique constraint,
and because the other services had (correctly) been made to wait for the
migration to \emph{succeed}, its failure took the whole stack down.

Demo data already existing is the normal state on a re-run, not an error:

\where{api/database/seeders/DatabaseSeeder.php}
\begin{lstlisting}[language=PHP]
if (Organization::query()->exists()) {
    $this->command?->info('Demo data is already present; nothing to seed.');

    return;
}
\end{lstlisting}

\note{Notice the interaction. Making dependencies stricter was correct, and it
converted a previously harmless failure into a total outage. Tightening one part
of a system can promote an existing minor fault into a major one --- which is an
argument for tightening things, not against it: the fault was always there.}

\chapter{Bugs that only appear when deploying}
\label{ch:bugs-deploy}

\section{The cookie that is never sent}

\textbf{Symptom (predicted, not suffered).} Deployed over plain HTTP, sign-in
appears to work and the session dies roughly fifteen minutes later. Nothing
errors. Nothing is logged.

\textbf{The cause.} The production configuration sets
\fpath{JWT\_REFRESH\_COOKIE\_SECURE=true}. A \fpath{Secure} cookie is only ever
sent back over HTTPS. Over plain HTTP the browser accepts it and then never
returns it --- so the first refresh, when the fifteen-minute access token
expires, has no credential.

Every component behaves exactly as specified. The API sets the cookie correctly.
The browser stores it correctly. It is simply never transmitted.

\warnbox{This is why the deployment guide treats TLS as mandatory rather than a
finishing touch. ``It works on HTTP, I will add HTTPS later'' produces an
application that appears to work and logs people out fifteen minutes later, in a
different component, with no error anywhere.}

\section{Three containers, three different signing keys}

\textbf{The setting.} Moving from Docker Compose to a platform that gives each
container its own filesystem.

\textbf{The cause.} In Compose, \fpath{api}, \fpath{horizon} and
\fpath{scheduler} share one volume holding the JWT key pair. On a platform
without shared volumes, each container generates its \emph{own} pair at startup.
A token signed by one then fails verification in another.

The symptom would be intermittent, inexplicable ``invalid signature'' errors
depending on which container answered.

\textbf{The fix.} Inject the key as an environment variable rather than
generating it. The configuration already supported this:

\where{api/config/jwt.php}
\begin{lstlisting}[language=PHP]
'private_key' => env('JWT_PRIVATE_KEY'),
'private_key_path' => env('JWT_PRIVATE_KEY_PATH', storage_path('keys/jwt-private.pem')),
\end{lstlisting}

\note{This is the shape of nearly every compose-to-platform port. The bug is not
in your code --- the code is correct. It is that a \emph{shared volume} is an
ambient guarantee Compose provides in one line and the platform silently removes.
The same is true of a flat IPv4 network and of \fpath{depends\_on} ordering. Each
has to become explicit configuration or it becomes a runtime mystery.}

\section{A hostname that resolves and still cannot be reached}

\textbf{The cause.} One platform's private DNS resolves service names to
\textbf{IPv6 only}. The API's nginx listened with \fpath{listen \$\{PORT\};},
which binds IPv4 alone. The hostname resolved perfectly; nothing could connect.

\textbf{The fix.} Bind both stacks.

\where{api/docker/render-nginx.conf.template}
\begin{lstlisting}[language=nginxconf]
listen ${PORT};
# Dual-stack. Railway's private DNS (*.railway.internal) resolves to IPv6
# only, so an IPv4-only listener is unreachable from a sibling service
# however correct the hostname is. Harmless on Render, required here.
listen [::]:${PORT};
\end{lstlisting}

\warnbox{Verifying this needed care. \fpath{nginx -t} only checks
\emph{syntax} --- a configuration that cannot bind still passes. The check that
matters is starting it and confirming both sockets are actually listening and
serving. A test that validates a file is not a test that validates a program.}

\section{A message that was a count, not a diagnosis}

\textbf{Symptom.} A container refused to start, reporting: \emph{database
unreachable after 30 attempts}.

That sentence contains no information. It cannot distinguish an unresolvable
hostname from a refused connection from a rejected password --- three completely
different problems with three completely different fixes. The probe had the real
error and threw it away:

\begin{lstlisting}[language=bash]
until php -r "new PDO(...);" 2>/dev/null; do
\end{lstlisting}

That \fpath{2>/dev/null} discarded the driver's message on every one of the
thirty attempts.

\textbf{The fix.} Keep the last error and report it.

\begin{lstlisting}[language=bash]
until db_error="$(php -r "new PDO(...);" 2>&1)"; do
  ...
  # Report what the driver actually said. "unreachable after 30 attempts" is
  # a count, not a diagnosis.
  echo "{\"level\":\"fatal\",\"message\":\"database unreachable\",\"last_error\":\"${detail}\"}" >&2
\end{lstlisting}

The two failure modes it had been conflating, now distinguishable:

\begin{lstlisting}
could not translate host name "nosuchhost" to address: Name does not resolve
connection to server at "postgres" failed: FATAL: password authentication failed
\end{lstlisting}

\note{Suppressing stderr in a retry loop is extremely common --- it keeps the
logs tidy while the loop retries. It also means that when the loop finally gives
up, the one moment you need the reason, it has been discarded thirty times over.
Keep the last error even when you discard the rest.}

\section{Two mistakes of mine, for completeness}

\subsection{The tool that reported success without running}

A formatting check reported ``passed'' every single time, locally --- because a
hook intercepted the command and returned a canned success. Every local ``lint
is clean'' was meaningless, and CI kept failing on formatting.

\warnbox{The lesson generalises well beyond this project: \textbf{a check that
cannot fail is not a check.} If you have never seen a tool report a failure, you
do not know it is running. Break something on purpose and confirm it complains.}

\subsection{The document that destroyed itself}

The first attempt at writing this book generated the LaTeX through a shell
heredoc. Every \fpath{\textbackslash\textbackslash} collapsed into a single
backslash. LaTeX builds table rows and line breaks out of
\fpath{\textbackslash\textbackslash}, so every table in the document was
silently destroyed --- and the error it produced pointed at a title line that
looked perfectly correct.

\note{Both of these share a shape with the screenshot bug in
Chapter~\ref{ch:bugs-frontend}: a tool reported success while producing wrong
output. That is the failure mode worth fearing most, because every instinct you
have says to trust it.}
